\documentclass[]{article}  

% Default font is the helvetica postscript font
\usepackage{helvet}
\usepackage{hyperref}
\usepackage{amsmath}
\usepackage{amsthm}
\usepackage{amssymb}
\usepackage[a4paper,margin=25mm]{geometry}
\usepackage{fancyhdr}
\usepackage{lipsum}
\usepackage[ruled,vlined]{algorithm2e}
\usepackage[utf8]{inputenc}
\usepackage[english]{babel}
\usepackage{graphicx}


\newtheorem{theorem}{Theorem}
\newtheorem{lemma}[theorem]{Lemma}
\newtheorem{defn}[theorem]{Definition}
\newtheorem{remark}[theorem]{Remark}
\newtheorem{corollary}[theorem]{Corollary}
%\renewcommand{\baselinestretch}{1.15}
\newcommand{\s}{\mathbf{s}}
\newcommand{\ba}{\mathbf{a}}
\newcommand{\bw}{\mathbf{w}}
\newcommand{\bx}{\mathbf{x}}
\newcommand{\bv}{\mathbf{v}}
\newcommand{\thetab}{\boldsymbol{{\theta}}}
\newcommand{\h}{\mathbf{h}}

\newcommand{\cA}{\mathcal{A}}
\newcommand{\cW}{\mathcal{W}}
\newcommand{\cH}{\mathcal{H}}

\newcommand{\R}{\mathbb{R}}
\newcommand{\E}{\mathbb{E}}
\newcommand{\Prob}{\mathbb{P}}
\newcommand{\KL}{\mathrm{D}_{KL}}



\title{Four Cute Facts on Basic Concentration Inequalities}
\author{Behrad Moniri\\\texttt{bemoniri@live.com}}
\date{}

\usepackage{sectsty} \sectionfont{\fontsize{15}{15}\selectfont}
\begin{document}
%\setcounter{section}{-1}
\maketitle

\begin{abstract}
Four cute fact on basic concentration inequalities are presented. First, we will provide a bound on the variance of a sub-Gaussian random variable. Then we go on to compare Chernoff and Moments bound, and also Bennett's and Bernstein's Inequalities. We finish this note by proving a bound on the expected value of a random variable satisfying a Bernstein's tail inequality. These results can be found (or are left as exercises) in well known textbooks such as \cite{concentration, wainwright2019high}.
\end{abstract}

\section{Bounding Variance of a  Sub-Gaussian Random Variable}

\begin{theorem}
Let $X$ be a zero-mean, $\sigma^2$-subGaussian random variable. $\mathrm{Var}[X] \leq \sigma^2$.
\end{theorem}

\begin{proof}
By the subGaussian property, for any $\lambda \in \R$ we have 
$\E[e^{\lambda Y}] \leq e^{\lambda^2 \sigma^2 / 2}$. We can write
\begin{align}
2 e^{\frac{\lambda^2\sigma^2}{2}} &\geq \E[e^{\lambda Y}] + \E[e^{-\lambda Y}]   \\
&=2 + 2\sum_{k = 1}^{\infty} \lambda^{2k} \frac{\E[Y^{2k}]}{(2k)!}\\
&\geq 2 + 2\lambda^2 \frac{\E[Y^2]}{2}.
\end{align}
Hence, for all $\lambda >0$, we have $ h(\lambda) = 1 + \frac{\lambda}{2} \E[Y^2] - e^{\lambda\sigma^2/2} \leq 0$. We have $h(0) = 0$, hence $h'(0)$ should be negative. Thus
\begin{align}
h'(\lambda) = \frac{1}{2}  \E[Y^2] - \frac{\sigma^2}{2}e^{\lambda \sigma^2/2}\Big|_{\lambda = 0} \leq 0 \implies \E[Y^2] \leq \sigma^2,
\end{align}
which concludes the proof.
\end{proof}

\begin{remark}
The equality $\E[Y^2] = \sigma^2$ might not hold even if $\sigma$ is the smallest number satisfying $\E[e^{\lambda Y}] \leq e^{\lambda^2 \sigma^2 / 2}$. To see this, let $X$ be the following discrete random variable
\begin{equation}
X = 
\begin{cases}
+1, & p\\
0, & 1-2p\\
-1, & p
\end{cases}
\end{equation}
It can easily be seen that $\E[X] = 0$ and $\E[X^2] = 2p$. Assume that the equality holds with $\sigma^2 = \mathrm{Var}[X] = 2p$.
\begin{align*}
g(\lambda) &\triangleq \log\E[e^{\lambda X}] - \frac{\lambda^2 \sigma^2}{2} - \lambda \mu\\
&= \log(pe^\lambda + p e^{-\lambda} + 1 - 2p) - \lambda^2 p \leq 0 \;\;\; \forall \lambda \in \R.
\end{align*}
Now, let  $p=0.1$ and $\lambda = 1$, then
$$\log(0.1e^1 + 0.1e^{-1} + 1 - 0.2) - 1 \time 0.1 \approx 0.00311 \geq 0$$
which is a contradiction.

\end{remark}


\section{Chernoff's vs. Moments Tail Bound}

\begin{theorem}
\label{cherrvsmom}
Suppose that $X\geq 0$ and that the moment generating function of $X$ in an interval around zero. Given $\delta>0$ and integer $k$, we have
\begin{equation}
\inf_{k=0, 1, 2, ...} \frac{\E[X^k]}{\delta^k} \leq \inf_{\lambda>0} \frac{\E[e^{\lambda X}]}{e^{\lambda \delta}}.
\end{equation}
\end{theorem}

\begin{proof}
We will prove that 
$\forall \lambda > 0: \inf_{k=0, 1, 2, ...} \frac{\E[X^k]}{\delta^k} \le 
\frac{\E[e^{\lambda X}]}{e^{\lambda \delta}}.$
The theorem follows from this statement. Define $y_k := \frac{\E[X^k]}{\delta^k}$ and $z_k = \frac{1}{k!} \lambda^k \delta^k$. By Taylor's expansion
\begin{equation}
\frac{\E[e^{\lambda X}]}{e^{\lambda \delta}} = 
\frac{\sum_{k=0}^{\infty} \frac{1}{k!} \lambda^k\E[X^k]}{
	\sum_{k=0}^{\infty} \frac{1}{k!} \lambda^k\delta^k
} = 
\frac{\sum_{k=0}^{\infty} \frac{1}{k!} \lambda^k\delta^k y_k}{
	\sum_{k=0}^{\infty} \frac{1}{k!} \lambda^k\delta^k
} = \frac{\sum_{k = 0}^{\infty} z_k y_k}{\sum_{k = 0}^{\infty}z_k} \geq \inf_{k=0,1,2,...} y_k,
\end{equation}
which proves the theorem.
\end{proof}

\begin{corollary}
We can write two different tail bounds for a positive variable $X$:
\begin{equation}
\Prob[X\geq t] \leq \inf_{k=0, 1, 2, ...}\frac{\mathbb{E}[X^k]}{t^k}
\end{equation}
and 
\begin{equation}
\Prob[X\geq t] \leq \inf_{\lambda>0}\frac{\E[e^{\lambda X}]}{e^{\lambda \delta}}.
\end{equation}
Theorem \eqref{cherrvsmom} states that the first bound is tighter than the second. However, the second bound is more popular, because of the ease of extending to to sum of i.i.d. random variables.
\end{corollary}




\section{Bennett's vs. Bernstein's Inequality}

	Let $X_i$ be a mean-zero random variables such that $|X_i| < b$.  First, We intend to bound the moment generating function from above.  For every $\lambda \in \R$ we can write

\begin{align*}
\E(e^{\lambda X_i}) &= 
1+	\E \Big(\sum_{k = 2}^{\infty}\frac{\lambda^k X_i^k}{k!}\Big)\\
&= 1 + \sum_{k = 2}^{\infty} \frac{\lambda^k \E(X_i^2 X_i^{k-2})}{k!}\\
&\leq 1 + \sum_{k = 2}^{\infty} \frac{\lambda^k \sigma_i^2 b^{k-2}}{k!}\\
&= 1 + \frac{\sigma_i^2}{b^2} \sum_{k = 2}^{\infty} \frac{\lambda^k b^k}{k!}\\
&= 1 + \frac{\sigma_i^2}{b^2} (e^{\lambda b} - \lambda b - 1)\\
&\leq \exp\Big(\frac{\sigma_i^2}{b^2} (e^{\lambda b} - \lambda b - 1)\Big).
\end{align*}
Hence 
\begin{equation}
\label{behrad:2.7.moment}
\log\E(e^{\lambda X_i}) \leq \frac{\sigma_i^2}{b^2} (e^{\lambda b} - \lambda b - 1).
\end{equation}


Since $X_i$s are independent random variables, we can write
\begin{align}{}
\label{behrad:2.7.tail}
\Prob\Big(\sum_{i = 1}^{n} X_i \geq t\Big) &\leq \frac{\E\Big(\exp(\sum_{i = 1}^{n} \lambda X_i)\Big)}{e^{\lambda t}} = \frac{\prod_{i = 1}^{n} \E \Big(\exp(\lambda X_i)\Big)}{e^{\lambda t}}
\end{align}

Substituting \eqref{behrad:2.7.moment} in \eqref{behrad:2.7.tail} we arrive at the following bound:

\begin{align*}
\mathbb{P} \Big(\sum_{i = 1}^{n} X_i \geq t \Big) &\leq e^{-\lambda t} \prod_{i = 1}^{n} \E(\exp(\lambda X_i))\\
&\leq e^{-\lambda t} \prod_{i = 1}^{n} \exp\Big(\frac{\sigma_i^2}{b^2}(e^{\lambda b} - 1 - \lambda b)\Big)
\end{align*}

Optimizing $\lambda$ yield the Bennett's inequality
\begin{equation}
\Prob\Big(\sum_{i = 1}^{n} X_i \geq t\Big) \leq \exp\Big(-\frac{n\sigma^2}{b^2} h\Big(\frac{bt}{n\sigma^2}\Big)\Big) ,
\end{equation}
where $h(t) = (1+t) \log(1+t)$ for $t>0$. Bernstein's inequality is a weakening of Bennett. Let $\Phi(x) = \frac{x^2}{2 + \frac{2}{3} x}$. It is easy to see that
\begin{equation}
\forall x \geq 0 \;\;\; (1+x) \log(1+x) - \frac{x^2}{2 + 2/3 x} \geq 0.
\end{equation}
Substituting $h$ with $\Phi$ in Bennett yield the Bernstein's inequality.

\section{Bernstein's Expectation}
\begin{theorem}
	Consider a random variable $Z$ that satisfies the following Bernstein-type inequality:
	\begin{equation}
	\Prob(Z \geq t) \leq C \exp\Big(\frac{-t^2}{2\sigma^2 + 2bt}\Big).
	\end{equation}
	We have 
	\begin{equation}
	\E[Z] \leq 4b(1 + \log(C)) + 2\sigma \big(\sqrt{\pi} + \sqrt{\log(C)}\big).
	\end{equation}
\end{theorem}

\begin{proof}
	If $\sigma^2 \leq bt$ then $\sigma^2 + bt \leq 2bt$. Otherwise $\sigma^2 + bt \leq 2\sigma^2$. We can rewrite the bound on $Z$ in the following form:
	\begin{align}
	\Prob(Z \geq t) &\leq C \exp\Big(\frac{-t^2}{2\sigma^2 + 2bt}\Big) \leq 				
	C\max\Big[\min\Big(\exp\big(\frac{-t^2}{4bt}\big),1\Big), \min\Big(\exp\big(\frac{-t^2}{4\sigma^2}\big),1\Big)	
	\Big]\\		
	&\leq C \Big[
	\min\Big(\exp\big(\frac{-t^2}{4bt}\big),1\Big) +  \min\Big(\exp\big(\frac{-t^2}{4\sigma^2}\big),1\Big)	
	\Big].
	\end{align}
	By writing the expected value as the integral of the tail probability function and bounding the tail probability, it is easy to derive a bound on $\E(Z)$.		
	\begin{align}
	\E[Z] &= \int_{0}^{\infty}\Prob(Z \geq t) \, dt\\ 
	&\leq  \int_{0}^{\infty}\min\Big\{1, C\exp\big(\frac{-t}{4b}\big)\Big\} \, dt +
	\int_{0}^{\infty} \min\Big\{1, C\exp\big(\frac{-t^2}{4\sigma^2}\big)\Big\} \, dt\\
	&= 4b(1 + \log(C)) + 2\sigma \big(\sqrt{\pi} + \sqrt{\log(C)}\big).
	\end{align}
	Where the last inequality follows from computing the integrals.
\end{proof}


\bibliography{bib}
\bibliographystyle{alpha}


\end{document}